\documentclass{article}
\usepackage{graphicx}
\usepackage{amsmath}
\usepackage[margin=2cm, top=2cm]{geometry}
\usepackage{tikz}
\usetikzlibrary{angles, quotes}
\usetikzlibrary{arrows.meta}
\usetikzlibrary{positioning}
\usepackage{pgfplots}
\pgfplotsset{compat=1.18}
\usepackage{amssymb}
\usepackage{enumitem}
\usepackage{tocloft}
\renewcommand{\cftsecdotsep}{\cftdotsep}
\cftpagenumbersoff{subsection}
\usepackage[normalem]{ulem}
\usepackage{dashundergaps}
\usepackage{xcolor}
\usepackage[most]{tcolorbox}
\usepackage{chemfig}
\usepackage[version=4]{mhchem}
\usepackage{siunitx}

\definecolor{defmain}{HTML}{1d1280}
\definecolor{defbg}{HTML}{F4F8F7}
\definecolor{rpmain}{HTML}{790a73}
\definecolor{rpbg}{HTML}{F5F5F8}
\definecolor{zgmain}{HTML}{146609}
\definecolor{zgbg}{HTML}{F4F6F8}

\newtcolorbox{defbox}[2][]{
    enhanced,
    colback=defbg,
    colframe=defmain,
    arc=3mm,
    boxrule=0.6pt,
    leftrule=1mm,
    before upper={\textbf{\sffamily\color{defmain}#2}\par\vspace{8pt}},
    #1
}

\newtcolorbox{rpbox}[2][]{
    enhanced,
    colback=rpbg,
    colframe=rpmain,
    arc=3mm,
    boxrule=0.6pt,
    leftrule=1mm,
    before upper={\textbf{\sffamily\color{rpmain}#2}\par\vspace{8pt}},
    #1
}

\newtcolorbox{zgledbox}[2][]{
    enhanced,
    colback=zgbg,
    colframe=zgmain,
    arc=3mm,
    boxrule=0.6pt,
    leftrule=1mm,
    before upper={\textbf{\sffamily\color{zgmain}#2}\par\vspace{8pt}},
    #1
}

\title{Zapiski za kemijo v 3. letniku srednje šole}
\date{}
\author{}

\begin{document}

\maketitle

% =====================================================================
\section{Ogljikovodiki}

\begin{defbox}{Tipi formul}
\ce{C3H8} \hspace{2.725cm} \textbf{Molekulska formula} \\
\ce{C3H8} \hspace{2.725cm} \textbf{Empirična}(Pokrajšana formula z največjim skupnik deljiteljem) \\
\ce{CH3-CH2-CH3} \hspace{1cm} \textbf{Racionalna formula} \\ \\
\chemfig{H-C(-[2]H)(-[6]H)-C(-[2]H)(-[6]H)-C(-[2]H)(-[6]H)-H} \hspace{2cm} \textbf{Strukturna formula} \\ \\ \\
\chemfig{-[:30]-[:-30]} \hspace{2cm} \textbf{Skeletna formula}
\end{defbox}

% TODO: stereo formula (klinasta vezava) -- v chemfig z < in <:
% \chemfig{C(-[:210]H)(<[:150]H)(<:[:330]H)(-[:30]H)}  % preveri

\begin{defbox}{Organske spojine}
Organske spojine vsebujejo predvsem elemente
\ce{C}, \ce{H}, \ce{O}, \ce{N}, \ce{P}, \ce{S} in halogene (npr.\ \ce{I}).
Najdemo jih v maščobah, mišicah \dots % TODO: dopolni
\end{defbox}

\subsection{Hibridizacija}

\begin{defbox}{Hibridizacija}
Hibridizacija je mešanje orbital ogljikovega atoma. % TODO: preveri besedilo v zvezku
\begin{center}
\begin{tabular}{lll}
\textbf{Hibridizacija} & \textbf{Vezi med atomi C} & \textbf{Skupina} \\ \hline
$sp^3$ & samo $\sigma$ vezi (enojna vez) & alkani \\
$sp^2$ & ena $\sigma$ in ena $\pi$ vez (dvojna vez) & alkeni \\
$sp$   & ena $\sigma$ in dve $\pi$ vezi (trojna vez) & alkini \\
\end{tabular}
\end{center}
\end{defbox}

% TODO: TikZ skica orbital (sp, sp2, sp3) in primer z oznakami hibridizacije
% nad vsakim C atomom (primer iz zvezka z dvojno in trojno vezjo)

\subsection{Razdelitev ogljikovodikov}

\begin{defbox}{Alifatski ogljikovodiki}
\begin{itemize}[noitemsep]
    \item \textbf{ciklični}
    \item \textbf{aciklični}
    \begin{itemize}[noitemsep]
        \item \textbf{nasičeni} -- alkani
        \item \textbf{nenasičeni} -- alkeni, alkini
    \end{itemize}
\end{itemize}
% TODO: aromatski ogljikovodiki (benzen, metilbenzen)
\end{defbox}

\begin{rpbox}{Splošne formule}
\begin{center}
\begin{tabular}{ll}
aciklični alkani: & $\mathrm{C}_n\mathrm{H}_{2n+2}$ \\
cikloalkani:      & $\mathrm{C}_n\mathrm{H}_{2n}$ \\
alkeni:           & $\mathrm{C}_n\mathrm{H}_{2n}$ \\
alkini:           & $\mathrm{C}_n\mathrm{H}_{2n-2}$ \\
\end{tabular}
\end{center}
\end{rpbox}

\begin{zgledbox}{Določitev formule iz molske mase}
Alkan ima $M_r = 72$. Koliko atomov C ima?
\begin{align*}
12n + 1\cdot(2n+2) &= 72 \\
14n + 2 &= 72 \\
n &= 5
\end{align*}
Formula: \ce{C5H12}.
\end{zgledbox}

\subsection{Alkani}

\begin{defbox}{Vrste ogljikovih atomov}
Glede na število vezanih ogljikovih atomov ločimo primarne ($1^\circ$),
sekundarne ($2^\circ$), terciarne ($3^\circ$) in kvarterne ($4^\circ$) atome C.
% TODO: primer s štetjem (npr. 6 x 1°, 2 x 2°, 2 x 3°, 1 x 4° pri 3-etil-2,2-dimetilpentanu)
\end{defbox}

\begin{rpbox}{Lastnosti alkanov}
\begin{itemize}[noitemsep]
    \item $\mathrm{C}_1$--$\mathrm{C}_4$: plin (metan, etan, propan, butan)
    \item $\mathrm{C}_5$--$\mathrm{C}_{16}$: tekočina
    \item $\mathrm{C}_{17}$ in več: trdna snov
    \item \textbf{nerazvejane} molekule $\Rightarrow$ višja vrelišča
    \item \textbf{razvejane} molekule $\Rightarrow$ nižja vrelišča
\end{itemize}
Sobni pogoji: $T = 20\,^\circ\mathrm{C}$, $p = 101{,}2\,\mathrm{kPa}$. % preveri vrednost tlaka
\end{rpbox}

\begin{defbox}{Nafta}
\begin{itemize}[noitemsep]
    \item rafiniranje nafte (daljše verige $\to$ krajše verige)
    \item frakcionirana destilacija nafte -- glej učbenik % TODO: skica destilacijskega stolpa
\end{itemize}
\end{defbox}

\begin{zgledbox}{Poskus: oktan in jod}
Oktan (bencin) in voda tvorita dve fazi (\emph{fazna meja}). Dodamo jod.
Ker so halogeni nepolarni, se v oktanu raztopi in oktanska faza se obarva vijolično.
% TODO: skica epruvete
\end{zgledbox}

% =====================================================================
\section{Izomerija}

\begin{defbox}{Izomerija}
\textbf{Izomerija} je pojav, ko imajo spojine enako molekulsko formulo, a različno
razporeditev atomov. Takšne spojine imenujemo \textbf{izomeri}.
Zaradi različne zgradbe se razlikujejo v fizikalnih in kemijskih lastnostih.
\end{defbox}

\subsection{Verižna izomerija}

\begin{defbox}{Verižna izomerija}
Pojavi se, ko je prisotna alkilna skupina (npr.\ \ce{CH3}).
\begin{center}
\begin{tabular}{cc}
\chemfig{-[:30](-[:90])-[:-30]-[:30]-[:-30]-[:30]} &
\chemfig{-[:30]-[:-30](-[:-90])-[:30]-[:-30]-[:30]} \\
2-metilheksan & 3-metilheksan \\
\end{tabular}
\end{center}
\end{defbox}

\begin{zgledbox}{Izomeri \ce{C6H14}}
\begin{center}
\begin{tabular}{ccc}
\chemfig{-[:30]-[:-30]-[:30]-[:-30]-[:30]} &
\chemfig{-[:30](-[:90])-[:-30]-[:30]-[:-30]} &
\chemfig{-[:30]-[:-30](-[:-90])-[:30]-[:-30]} \\
heksan & 2-metilpentan & 3-metilpentan \\[1.5em]
\chemfig{-[:30](-[:90])-[:-30](-[:-90])-[:30]} &
\chemfig{-[:30](-[:90])(-[:150])-[:-30]-[:30]} & \\
2,3-dimetilbutan & 2,2-dimetilbutan & \\
\end{tabular}
\end{center}
Vseh pet je \textbf{verižnih izomerov}.
% TODO: naloga z razvrstitvijo po vrelišču (3-metilheptan, oktan, 2,2,4-trimetilpentan)
\end{zgledbox}

\subsection{Položajna izomerija}

\begin{defbox}{Položajna izomerija}
Pojavi se, ko so prisotne druge skupine (npr.\ halogeni), ter ko se premaknejo
dvojne ali trojne (večkratne) vezi.
\begin{center}
\begin{tabular}{cc}
\chemfig{Cl-[:-30]-[:30]-[:-30]-[:30]} &
\chemfig{-[:30](-[:90]Cl)-[:-30]-[:30]} \\
1-klorobutan & 2-klorobutan \\[1.5em]
\ce{CH2=CH-CH2-CH3} & \ce{CH3-CH=CH-CH3} \\
but-1-en & but-2-en \\
\end{tabular}
\end{center}
% TODO: 1-, 2-, 3-fluoropentan
\end{defbox}

\subsection{Funkcionalna izomerija}

\begin{defbox}{Funkcionalna izomerija}
Pojavi se, ko se skupine zamenjajo. % TODO: dopolni z zgledi, ko jih obravnavate
\end{defbox}

% =====================================================================
\section{Gorenje ogljikovodikov}

\begin{defbox}{Vrste gorenja}
\begin{center}
\begin{tabular}{lll}
\textbf{popolno gorenje} & produkta: & \ce{CO2} in \ce{H2O} \\
\textbf{nepopolno gorenje} & produkta: & \ce{CO} in saje (\ce{C}) \\
\end{tabular}
\end{center}
\end{defbox}

\begin{rpbox}{Urejanje enačb}
Enačbo urejamo po vrstnem redu: 1) \ce{C}, 2) \ce{Cl} (če obstaja), 3) \ce{H}, 4) \ce{O}.
\end{rpbox}

\begin{zgledbox}{Popolno gorenje}
\begin{align*}
&\ce{2C4H10(g) + 13O2(g) -> 8CO2(g) + 10H2O(g)} \tag{butan} \\
&\ce{2C8H18(g) + 25O2(g) -> 16CO2(g) + 18H2O(g)} \tag{oktan}
\end{align*}
\end{zgledbox}

\begin{zgledbox}{Nepopolno gorenje}
Primer: toluen (metilbenzen) $+\ \ce{O2} \to \ce{CO}$ + saje. % TODO: uredi enačbo
\begin{center}
\chemfig{*6(-(-CH_3)=-=-=)}
\end{center}
\ce{CO} je strupen, ker se veže na hemoglobin v krvi in s tem prepreči vezavo kisika.
Je brez barve in vonja.
% TODO: razmerje ogljika/vodika proti kisiku -- plamen s sajami (preveri zapis v zvezku)
\end{zgledbox}

\begin{zgledbox}{Naloga}
Pri gorenju 1 mol nekega ogljikovodika nastane 5 mol \ce{CO2} in 4 mol \ce{H2O}.
\begin{align*}
&\mathrm{C}_x\mathrm{H}_y + \mathrm{O}_2 \to 5\,\mathrm{CO}_2 + 4\,\mathrm{H}_2\mathrm{O} \\
&\ce{C5H8 + 7O2 -> 5CO2 + 4H2O}
\end{align*}
Ogljikovodik je \ce{C5H8} (npr.\ pent-1,2-dien).
\end{zgledbox}

% =====================================================================
\section{Radikalska substitucija}

\begin{defbox}{Reakcije ogljikovodikov}
Alkani reagirajo na dva načina: \textbf{gorenje} in \textbf{radikalska substitucija}.
Pogoji za radikalsko substitucijo: povišana temperatura ($\Delta$) ali svetloba ($h\nu$).
\end{defbox}

\subsection{Prekinitev vezi}

\begin{defbox}{Vrste prekinitev vezi}
\textbf{Homolitska prekinitev}: nastaneta dva radikala.
\[ \ce{Br-Br} \longrightarrow \ce{Br^.} + \ce{^.Br} \]
\textbf{Heterolitska prekinitev}: nastaneta kation in anion.
\[ \ce{Br-Br} \longrightarrow \ce{Br+} + \ce{Br-} \]
Radikali nastanejo s segrevanjem ali svetlobo.
Energija svetlobe: $E = h\cdot\nu$, kjer je $h = 6{,}626\cdot 10^{-34}\,\mathrm{J\,s}$
(Planckova konstanta).
\end{defbox}

\begin{defbox}{Nukleofil in elektrofil}
\textbf{Nukleofil} je delec, ki išče pozitivni center (polarne molekule, npr.\ \ce{H2O}). \\
\textbf{Elektrofil} je delec, ki išče negativni center in sprejema elektrone (nepolarne molekule). \\
Molekule so lahko tudi brez naboja.
\end{defbox}

\subsection{Stopnje reakcije}

\begin{rpbox}{Stopnje radikalske substitucije}
\begin{enumerate}[noitemsep]
    \item \textbf{iniciacija} -- začetek
    \item \textbf{propagacija} -- potek (traja)
    \item \textbf{terminacija} -- konec
\end{enumerate}
\end{rpbox}

\begin{defbox}{Substrat in reagent}
\textbf{Substrat} je organska spojina, na katero delujemo z reagentom. \\
\textbf{Reagent} zamenja del substrata; je lahko element ali spojina. \\
\textbf{Reakcijski pogoji} so okoliščine, pri katerih poteka kemijska reakcija.
\end{defbox}

\subsection{Primeri}

\begin{zgledbox}{Kloriranje metana}
Kloriramo lahko enkrat ali večkrat.
\begin{align*}
\ce{CH4 + Cl2} &\xrightarrow{\ \Delta\ \text{ali}\ h\nu\ } \ce{CH3Cl + HCl} \quad \text{(klorometan)} \\
\ce{CH3Cl + Cl2} &\xrightarrow{\ \Delta\ \text{ali}\ h\nu\ } \ce{CH2Cl2 + HCl} \quad \text{(diklorometan)}
\end{align*}
\end{zgledbox}

\begin{zgledbox}{Kloriranje kloroetana}
\[ \ce{CH3-CH2Cl + Cl2} \xrightarrow{\ \Delta\ } \ce{CH3-CHCl2} + \ce{CH2Cl-CH2Cl} + \ce{HCl} \]
Nastaneta 1,1-dikloroetan in 1,2-dikloroetan.
\end{zgledbox}

\begin{zgledbox}{Kloriranje cikloheksana}
\begin{center}
\schemestart
\chemfig{*6(------)}
\arrow{->[\ce{Cl2}, $\Delta$][\ce{-HCl}]}
\chemfig{*6(--(-Cl)----)}
\schemestop
\end{center}
cikloheksan $\to$ klorocikloheksan
\end{zgledbox}

\begin{zgledbox}{Bromiranje metilbenzena}
\begin{center}
\schemestart
\chemfig{*6(-(-CH_3)=-=-=)}
\arrow{->[\ce{Br2}, $h\nu$][\ce{-HBr}]}
\chemfig{*6(-(-CH_2Br)=-=-=)}
\schemestop
\end{center}
Metilna skupina je alkan, benzenski obroč je aromatska spojina (ostane nespremenjen).
Produkt: (bromometil)benzen.
\end{zgledbox}

% =====================================================================
\section{Alkeni in alkini}

\begin{rpbox}{Splošne formule}
\begin{center}
\begin{tabular}{ll}
alkeni: & $\mathrm{C}_n\mathrm{H}_{2n}$ \\
alkini: & $\mathrm{C}_n\mathrm{H}_{2n-2}$ \\
\end{tabular}
\end{center}
\end{rpbox}

\begin{rpbox}{Poimenovanje}
% TODO: dopolni po učiteljevih pravilih
Končnica \textbf{-en} (dvojna vez), \textbf{-in} (trojna vez), \textbf{-dien}, \textbf{-trien} \dots \\
Verigo oštevilčimo tako, da ima prva večkratna vez čim manjšo številko.
\end{rpbox}

\begin{zgledbox}{Primeri}
\begin{center}
\begin{tabular}{cc}
\chemfig{CH_3-CH(-[2]CH_3)-CH_2-C~CH} &
\chemfig{CH_3-C(-[2]CH_3)=CH-CH=CH_2} \\
4-metilpent-1-in & 4-metilpenta-1,3-dien \\
\end{tabular}
\end{center}
% TODO: ostale naloge iz zvezka (2,4,6-trimetilhepta-1,3,5-trien,
% 3-metilciklobut-1-en, 1,4-dimetilciklobut-1-en, 2,5-dimetilcikloheksa-1,3-dien,
% 2,6,6-trimetilhept-3-in, 7-metilnona-3,5-dien)
\end{zgledbox}

\end{document}